
\begin{proposition}[Refuted]\label{cdp:st:D.1.6}

\textbf{A second count counterexample under the integral hypotheses.}
Wang's attached construction $X_2(\lambda)$ has integral homology
$H_0=H_6=\Z$, $H_2=H_3=\Z/2\Z$, and the other reduced
homology groups zero. Hence its integral cohomology satisfies
$H^1=H^2=0$. Its fibration has four exceptional values, over which
the reduced fibres are irreducible: the multiplicity-three
bielliptic fibre, two rank-one fillings, and the rank-two filling.
In the rank-one filling over the $I_2$ value, the primitive period
exponent is $(2,1)$; the component count is one, not two.
These are the construction's statements
\texttt{intro:main}, \texttt{fill:rankone},
\texttt{fill:a2}, \texttt{fill:log}, and
\texttt{main:theorem} in the attached primary source
\cite{Wang2026}.

In its cohomology basis the four displayed monodromies are
\[
\begin{aligned}
B_*&=\begin{pmatrix}1&0&-6&-2\\0&-1&-1&1\\0&1&0&-1\\0&0&0&1\end{pmatrix},&\qquad
B_2&=\begin{pmatrix}1&0&6&3\\0&1&2&1\\0&0&1&0\\0&0&0&1\end{pmatrix},\\
B_+&=\begin{pmatrix}1&0&0&0\\0&2&1&0\\0&-1&0&0\\0&0&0&1\end{pmatrix},&\qquad
B_-&=\begin{pmatrix}1&0&0&-1\\0&1&1&0\\0&0&1&0\\0&0&0&1\end{pmatrix}.
\end{aligned}
\]
Their product is one and their common invariant cohomology group
has rank one. On homology the matrices are $B_i^{-t}$; the Smith
diagonal of their difference matrix is $(1,1,1,0)$. Consequently
$t'=1$ and the corrected count is $r=4-1+1=4$, while the printed
count gives seven. The exact matrix calculation is included in the
program. The topological construction and its local fillings are
being used here as cited primary-source results, not inferred from
the matrices alone. Its coherent direct images are derived separately
in Statement \ref{cdp:st:D.1.13} below.
\end{proposition}


